\section{Conclusion and Future Work}
\label{sec:conclusion}

We have presented \alaya{}, a memory engine for conversational AI agents
that treats memory as a dynamic process rather than a static store. Drawing
on complementary learning systems theory, Bjork dual-strength forgetting,
Hebbian learning, and Yogac\=ara Buddhist psychology, \alaya{} implements
five lifecycle operations---consolidation, transformation, forgetting,
preference emergence, and emergent ontology---that manage memory health
over extended interaction horizons.

The emergent ontology contribution is particularly novel: categories arise
from dual-signal clustering (embedding similarity + graph topology) during
the transformation lifecycle, without supervision or manual curation.
Categories participate in the Hebbian graph, enabling spreading activation
to surface categorically related memories during retrieval.

Our baseline evaluation reveals a statistically significant retrieval
crossover between full-context injection and na\"ive RAG---neither extreme
is universally optimal---motivating the lifecycle-aware memory management
that \alaya{} provides. The system is implemented as a zero-dependency Rust
library requiring no external databases or network services, available as
an open-source crate at \url{https://crates.io/crates/alaya}.

\subsection*{Future Work}

\textbf{v0.2.0} will extend the emergent ontology with Yogac\=ara-grounded
enhancements: category hierarchy (\emph{vikalpa}---conceptual
construction), prototype-based membership with graded typicality
(\emph{n\=ama-r\=upa}---name and form), category seeds that retain dormant
state and can re-emerge (\emph{b\=ija}---seeds), and conceptual
transformation where categories themselves evolve through centroid drift
(\emph{\=a\'sraya-par\=av\d{r}tti}). Cross-domain bridging through
category nodes will enable analogical reasoning---``cooking techniques''
surfacing chemistry-related memories through shared category links.

\textbf{End-to-end evaluation} of \alaya{}'s full retrieval pipeline
(BM25 + vector + graph + Hebbian reranking) with lifecycle operations
(consolidation, forgetting, category-boosted retrieval) against LoCoMo
and LongMemEval baselines is the immediate empirical priority.

\textbf{Learned thresholds} for category assignment, discovery, merge, and
dissolution---replacing the current hand-tuned defaults with values
optimized from task performance---represent a natural extension of the
RL-trained memory management trend identified in the companion
survey~\cite{hui2026taxonomy}.
